Multi-agent LLM systems increasingly power critical reasoning tasks, yet most deployments use
homogeneous model ensembles that may share correlated failure modes.
Classical ensemble theory predicts that agent diversity reduces correlated errors,
but whether \emph{capability-level} diversity within a single model family provides
measurable benefits remains an open empirical question.
We study this question by constructing \advbench, a 30-question adversarial reasoning benchmark
spanning five cognitive trap categories---misleading math, causal traps,
logical deception, numerical tricks, and framing effects---and evaluating
five experimental conditions using real API calls to \haiku and \sonnet models.
Our conditions cover a diversity continuum from single agents, to homogeneous three-agent ensembles
using majority voting, to heterogeneous ensembles mixing small and large models,
and a structured debate protocol.
We find that both individual models achieve high accuracy (93.3\%) on our benchmark,
leaving little room for diversity-based improvement.
Error correlation between the cross-capability \haiku--\sonnet pair
($r{=}0.464$) is measurably lower than within-capability pairs ($r{=}0.695$),
partially supporting the diversity hypothesis.
However, heterogeneous ensembles match---rather than exceed---homogeneous ensembles
at 96.7\% accuracy ($p{=}1.000$, McNemar), and the debate protocol
slightly degrades performance due to sycophantic revision.
The dominant shared failure mode is a safety-aligned refusal to recommend
a surgery without full clinical context, shared identically across
all five conditions.
Our results suggest that within-family capability diversity alone is insufficient
to overcome shared alignment properties, and that truly complementary failure modes
require cross-family architectural diversity.
